\section{Exact observation and optimal local placement}
\label{sec:exact-observation}
Knowing the offset changes the design problem: mobility can now be
placed around the observed kinetic zeros. We first solve the local
quadratic problem over all integrable coefficients, then use explicit
policies to justify averaging over offsets. The local model is an
energy-dual problem on the line, not a stationary transport process on
an infinite wall. Saturated-gradient certificates are an established
method in reinforcement design \citep{ButtazzoOudetVelichkov2015}.
Constant optimal gradients also occur in fixed-volume cooling-fin
design \citep{AlexandersenSigmund2021}. The profile and constants
needed here are derived directly for quadratic killing and a uniform
source.

\subsection{The quadratic placement problem and its global certificate}
For $a,m>0$, define
\begin{equation}\label{eq:placement-line}
 \mathscr J_a^{\rm quad}(D)=\sup_{v\in C_c^\infty(\R)}
 \left\{2\int_\R v\dd x-\int_\R
                 [D|v'|^2+ax^2v^2]\dd x\right\},
 \qquad D\ge0,\quad\int_\R D\dd x=m.
\end{equation}
As in \eqref{eq:J-definition}, an infinite response is allowed, and
no operator realization for an arbitrary $L^1$ coefficient is presumed.

\begin{theorem}[Exact quadratic placement]\label{thm:quadratic-placement}
Set
\begin{equation}\label{eq:placement-radius}
 R=\left(\frac{80m}{3a}\right)^{1/5},\qquad
 y=\frac{|x|}{R},\qquad C_{\rm pl}=12\left(\frac3{80}\right)^{1/5}.
\end{equation}
The unique minimizer of \eqref{eq:placement-line}, up to equality almost
everywhere, is
\begin{equation}\label{eq:placement-profile}
 D_*(x)=\frac{aR^4}{8}\,y(1-y)^2(1+2y)\ind_{\{y<1\}}.
\end{equation}
Its source representative and optimal response are
\begin{align}
 h_*(x)&=
 \begin{cases}
 (9-8y)/(aR^2),&y\le1,\\
 1/(ax^2),&y>1,
 \end{cases}\label{eq:placement-field}\\
 \inf_D\mathscr J_a^{\rm quad}(D)
 &=\mathscr J_a^{\rm quad}(D_*)
 =\frac{12}{aR}=C_{\rm pl}a^{-4/5}m^{-1/5}.
 \label{eq:placement-value}
\end{align}
The source representative is understood in the completion in the
reaction-plus-derivative energy. In particular, the theorem does not
regard the constant source as an $L^2(\R)$ vector.
\end{theorem}
\begin{proof}
Write $p(y)=y-3y^3+2y^4$. It is nonnegative on $[0,1]$, and
\begin{equation}\label{eq:placement-identities}
 \int_0^1p(y)\dd y=\frac3{20},\qquad
 p'(y)+y^2(9-8y)=1.
\end{equation}
The first identity gives $\int D_*=3aR^5/80=m$. On $0<x<R$,
$D_*h_*'=-Rp(x/R)$, so the second proves
$-(D_*h_*')'+ax^2h_*=1$. Reflection gives the same equation on
$(-R,0)$; outside $[-R,R]$ it reduces to $ax^2h_*=1$.
The flux vanishes at $-R,0,R$. Thus neither the cusp at zero nor the
change of slope at the support edges introduces a delta source.
Direct integration yields
\begin{equation}\label{eq:placement-energies}
 \int h_* =\frac{12}{aR},\qquad
 \int D_*|h_*'|^2=\frac{12}{5aR},\qquad
 \int ax^2h_*^2=\frac{48}{5aR}.
\end{equation}
The inside and outside contributions to $\int h_*$ are $10/(aR)$
and $2/(aR)$, respectively.

For every compact smooth $v$, integration of the weak equation gives
\[
 2\int v-\int[D_*|v'|^2+ax^2v^2]
 =\int h_*-\int[D_*|(v-h_*)'|^2+ax^2(v-h_*)^2]
 \le\int h_*.
\]
The field $h_*$ belongs to the stated energy completion: cut off its
$|x|^{-2}$ tail at distance $L\to\infty$, then mollify its finitely
many corners. Tail source and reaction integrals tend to zero, and
the derivative energy is supported in $[-R,R]$, where $D_*$ is bounded.
These approximations approach $h_*$ in energy and in source integral,
proving equality in \eqref{eq:placement-value} for $D_*$.

The bound that excludes every competing placement is
\begin{equation}\label{eq:placement-slope}
 b_*:=\|h_*'\|_\infty=\frac8{aR^3},\qquad
 |h_*'|=b_*\ \ (0<|x|<R),\qquad
 |h_*'|<b_*\ \ (|x|>R).
\end{equation}
For an arbitrary nonnegative $D\in L^1$ of mass $m$, the same
cutoffs and mollifications have derivatives bounded by $b_*+o(1)$
and converging almost everywhere to $h_*'$. Dominated convergence
against $D\dd x$ therefore makes them legitimate tests and gives
\begin{align}
 \mathscr J_a^{\rm quad}(D)
 &\ge2\int h_*-\int ax^2h_*^2-\int D|h_*'|^2\notag\\
 &\ge2\int h_*-\int ax^2h_*^2-mb_*^2
 =\frac{12}{aR}.
 \label{eq:placement-certificate}
\end{align}
The derivative contribution at the large cutoff is $O(L^{-3})$;
mollification does not increase its Lipschitz bound. This justifies
\eqref{eq:placement-certificate} even for unbounded or highly
concentrated competitors.

If equality holds, the strict inequality in
\eqref{eq:placement-slope} forces $D=0$ almost everywhere outside
$[-R,R]$. Also $h_*$ maximizes the extended quadratic functional for
$D$: its value equals the supremum, and $h_*+t\phi$ is a legitimate
approximated test for every $\phi\in C_c^\infty$ and real $t$.
The first variation gives
$-(Dh_*')'+ax^2h_*=1$ in distributions. Hence $Dh_*'$ has an
absolutely continuous representative locally, with zero exterior
flux. On each open half-support, $h_*'$ is a nonzero constant;
integration from $R$ or $-R$ uniquely recovers
\eqref{eq:placement-profile}. This proves uniqueness almost everywhere.
As a check on the mass normalization,
$\partial_m\inf\mathscr J_a^{\rm quad}=-64/(a^2R^6)=-b_*^2$.
\end{proof}

The optimal coefficient vanishes linearly at the kinetic zero and
quadratically at its outer edges. These zeros do not invalidate the
certificate, which only uses smooth tests and zero flux. They do mean
that positivity or transmission across the center must not be imposed
silently as an extra constraint. A fixed positive mobility floor,
a pointwise cap, or a fixed fabrication length changes the problem.

\subsection{Separated defects on a compact wall}
The next result uses actual local masses of arbitrary competing fields.
It therefore does not assume beforehand that an optimal sequence
concentrates near the defects.
\begin{theorem}[Compact-wall placement and allocation]
\label{thm:compact-placement}
Let $k$ be continuous on the compact wall $\Gamma$, positive except at
finitely many distinct points $s_1,\ldots,s_n$, with $n\ge1$, and suppose
\begin{equation}\label{eq:placement-local-rate}
 k(s_j+x)=a_jx^2[1+o(1)],\qquad a_j>0,\qquad x\to0.
\end{equation}
Then, over all nonnegative $L^1(\Gamma)$ fields of mass $M$,
\begin{equation}\label{eq:compact-placement}
 \inf_D J_k(D)\sim
 C_{\rm pl}\left(\sum_{j=1}^na_j^{-2/3}\right)^{6/5}M^{-1/5}.
\end{equation}
A matching leading allocation to the individual defects is
\begin{equation}\label{eq:placement-allocation}
 m_j=M\frac{a_j^{-2/3}}{\sum_i a_i^{-2/3}}.
\end{equation}
The kinetic profile and the separated defect positions are fixed as
$M\downarrow0$. The theorem does not apply uniformly as defects merge.
\end{theorem}
\begin{proof}
Fix $0<\eta<1/2$. Choose disjoint fixed neighborhoods
$U_j=\{s_j+x:|x|<r_j\}$ in which
$a_{j,-}x^2\le k(s_j+x)\le a_{j,+}x^2$, where
$a_{j,\pm}=(1\pm\eta)a_j$. For an arbitrary competitor let
$b_j=\int_{U_j}D$, so $\sum_jb_j\le M$.
For $b_j>0$, take the field \eqref{eq:placement-field} with parameters
$(a_{j,+},b_j)$ and multiply it by a fixed smooth cutoff equal to one
on $|x|<r_j/2$ and supported in $U_j$.
For all small $M$, uniformly over $0<b_j\le M$, the cutoff acts
only on the tail $1/(a_{j,+}x^2)$. It changes the source-minus-reaction
integral by $O_{\eta}(1)$ and introduces only bounded slopes, whereas
the central slope diverges at least as $M^{-3/5}$.
Mollifying these compact Lipschitz tests retains that slope bound.
Potential ordering and the actual mass $b_j$ in each neighborhood give
\begin{equation}\label{eq:placement-local-lower}
 J_k(D)\ge
 \sum_{j=1}^n C_{\rm pl}a_{j,+}^{-4/5}b_j^{-1/5}-C_\eta.
\end{equation}
The tests have disjoint supports, so their source and energy terms
add. If $b_j=0$, use the same test with an auxiliary positive budget
tending to zero: its mobility cost is zero and its
source-minus-reaction value diverges. Thus \eqref{eq:placement-local-lower}
holds with the value $+\infty$ in this case.

For positive weights $w_j$, elementary minimization gives
\begin{equation}\label{eq:placement-mass-minimization}
 \inf_{b_j>0,\,\sum b_j\le M}\sum_j w_jb_j^{-1/5}
 =M^{-1/5}\left(\sum_jw_j^{5/6}\right)^{6/5}.
\end{equation}
The objective decreases in each mass, so the total mass is $M$;
the derivative condition gives $b_j\propto w_j^{5/6}$, and strict
convexity proves the minimum. Substitution into
\eqref{eq:placement-local-lower} yields the unrestricted lower bound.

For the upper bound allocate the masses \eqref{eq:placement-allocation}
and insert \eqref{eq:placement-profile}, with curvature $a_{j,-}$
and mass $m_j$, into each neighborhood; set $D=0$ elsewhere.
The supports $|x|<R_j$ are disjoint for small $M$ and the budget
is exactly $M$. On a support, denote the interior field in
\eqref{eq:placement-field} by $h_{j,-}$. Its zero-flux equation and
$k\ge a_{j,-}x^2$ imply, for every smooth test restricted to that support,
\begin{equation}\label{eq:placement-support-upper}
 2\int_{-R_j}^{R_j}v-
 \int_{-R_j}^{R_j}[D|v'|^2+kv^2]
 \le\int_{-R_j}^{R_j}h_{j,-}=\frac{10}{a_{j,-}R_j}.
\end{equation}
This is direct completion of the square; no boundary trace of $v$
is constrained. Outside the supports use $2v-kv^2\le1/k$.
In $U_j\setminus[-R_j,R_j]$, the reciprocal integral is at most
$2/(a_{j,-}R_j)+O_\eta(1)$, and the remaining fixed region contributes
a bounded amount. Summing proves the upper counterpart of
\eqref{eq:placement-local-lower}. Divide by $M^{-1/5}$, first send
$M\to0$, then $\eta\to0$ to prove \eqref{eq:compact-placement}.
The bounded errors at fixed $\eta$ do not imply an additive bounded
error after subtracting the final leading term with the exact $a_j$.
\end{proof}

For comparison, the potential-ordering and interval-bracketing proof
of \cref{prop:harmonic} uses only the local quadratic asymptotic and
a positive lower bound away from the zeros. It therefore applies to
the present continuous rates and gives
$J_k(M/P)\sim C_0P^{1/4}\sum_ja_j^{-3/4}M^{-1/4}$.
Placement changes the divergence exponent from $1/4$ to $1/5$ for a
fixed set of observed quadratic defects.

\begin{corollary}[Physical localized trials and their bulk correction]
\label{cor:placement-bulk}
Under the physical assumptions of \cref{sec:model} with $V\ne0$,
the infimum of $D_\flow$ over physical budget-$M$ fields has leading
value $\chi$ times \eqref{eq:compact-placement}.
For the fixed-$\eta$ localized trials used in its proof, the minimal
closed derivative form is a valid realization and
\begin{equation}\label{eq:placement-bulk-remainder}
 \|kh-1\|_{L^2(\Gamma)}=O_\eta(M^{1/10}),\qquad
 \mathcal R(D,k)\longrightarrow\mathcal R_0,
\end{equation}
where $\mathcal R_0$ is defined by \eqref{eq:regular-bulk-problem}.
These are fixed-kinetic-profile statements, not bounds uniform in an
offset approaching a fold.
\end{corollary}
\begin{proof}
The leading full-bulk optimum follows immediately from
\cref{thm:optimal-transfer} for a single fixed realization and
\eqref{eq:compact-placement}; its positive-background mixture
preserves the budget and asymptotic value.
For the stronger assertion about the original degenerate trials,
closability can be checked on every compact subinterval where $D>0$.
If smooth $v_l\to0$ in $L^2$ and $\sqrt D\,v_l'$ is Cauchy in $L^2$,
its limit vanishes on each such subinterval by distributional
differentiation. It is zero where $D=0$, and the finitely many boundary
points have measure zero. Thus the derivative preform is closable.

The endpoints of each half-support have zero derivative capacity:
a transition of height one at an outer quadratic zero over width
$\epsilon$ costs $O(\epsilon)$; near the central linear zero, a
transition with derivative proportional to $1/x$ on
$\epsilon^2<x<\epsilon$ costs $O(1/|\log\epsilon|)$.
The corresponding $L^2$ costs tend to zero. Smooth approximations to
these transitions first localize bounded form functions to individual
components: the cutoff derivative cost is bounded by the squared
supremum of the function times the transition cost, while its remaining
energy on the shrinking transition intervals tends to zero.
Truncations $T_N(v)=\max(-N,\min(v,N))$ converge to $v$ in the form
norm as $N\to\infty$, extending this conclusion to all form-domain
elements. Thus the minimal closure permits independent restrictions
to the half-supports and exterior, with no endpoint matching condition;
finite endpoint traces need not exist. This domain fact is used only
for the explicit coefficients, not for arbitrary competing weights.

The reaction-plus-derivative form is coercive for each fixed $M>0$.
Indeed, away from the centers $k$ is bounded below on the relevant
closed sets. Near a center, anchor a test in an interval inside the
half-support where $D$ and $k$ are positive. Weighted Cauchy--Schwarz
and $D(x)\asymp_M |x|$ give
$|v(x)|^2\le C_M\mathfrak a_{k,D}[v](1+|\log |x||)$;
integrating this bound controls its remaining $L^2$ norm.
Thus its inverse $h$ exists in $L^2$ and in the minimal form domain.
Positivity and potential comparison on each support give
$0\le h\le h_{j,-}$: subtract the two weak equations and test with
the positive part of $h-h_{j,-}$, using the zero-capacity endpoints.
Outside the supports $h=1/k$ almost everywhere. Since
\[
 \max_{0\le y\le1}y^2(9-8y)=\frac{27}{16},
 \qquad
 0\le kh\le\frac{a_{j,+}}{a_{j,-}}\frac{27}{16}
 \quad\hbox{on support }j,
\]
$kh-1$ is bounded and supported on a set of length $O_\eta(M^{1/5})$.
This proves the first assertion in \eqref{eq:placement-bulk-remainder}.
No transmission rule for an initial point mass at a kinetic zero is
needed for this stationary, absolutely continuous formulation.

In \eqref{eq:bulk-remainder}, the load $\mathcal L_D$ consequently
converges in the mean-zero bulk energy dual norm to
$\mathcal L_0(f)=\int_\Omega(u-V)f-KV\int_\Gamma f_\Gamma$.
Dropping the nonnegative Schur penalty gives
$\limsup\mathcal R\le\mathcal R_0$.
For any fixed smooth bulk test $f$, choosing $v=f_\Gamma$ in that
penalty gives
$0\le\mathcal S_D(f_\Gamma)\le M\|\partial_sf_\Gamma\|_\infty^2\to0$.
Insert such tests in the supremum and use their density in the
mean-zero bulk energy space to get $\liminf\mathcal R\ge\mathcal R_0$.
This proves convergence without assuming a bounded tangential
derivative for the bulk maximizer.
\end{proof}

\subsection{The exact-observation mean and a measurable recovery policy}
Return to the cosine ensemble \eqref{eq:cosine-ensemble} and define
\begin{equation}\label{eq:oracle-value}
 \mathcal O(M)=F_1(M;c)
 =\inf_{c\mapsto D_c\ \mathrm{Borel},\ D_c\ge0,\ \int D_c=M}
      \frac14\int_{-2}^2J_c(D_c)\dd c.
\end{equation}
The budget is the same for every observed offset. This definition is
an infimum over policies, rather than an unproved interchange of
expectation and pointwise minimization.

\begin{theorem}[Value of exact kinetic observation]\label{thm:oracle}
As $M\downarrow0$,
\begin{equation}\label{eq:oracle-sharp}
 \mathcal O(M)\sim K_{\rm obs}M^{-1/5},\qquad
 K_{\rm obs}=\frac{2^{6/5}C_{\rm pl}}4
                 \Bfun\!\left(\frac12,\frac15\right).
\end{equation}
The full-bulk optimum is
$G_1(M;c)\sim\chi K_{\rm obs}M^{-1/5}$ in the convention of
\cref{thm:optimal-transfer}. Compared with the predetermined mean,
\begin{equation}\label{eq:oracle-blind-ratio}
 \frac{\mathcal O(M)}{\mathcal P_1(M)}
 \sim\frac{K_{\rm obs}}{K_1}M^{1/20}
 =1.21856579\ldots\,M^{1/20},
\end{equation}
where $K_{\rm obs}=22.40462823\ldots$ and $K_1=18.38606365\ldots$.
The same ratio holds for the two full-bulk mean optima.
\end{theorem}

We give the policy construction and its uniform estimates before
completing the proof. Let $t=1-|c|$ and fix a sufficiently large number
$L$. All constants below may depend on this fixed $L$, but not on $M$
or $c$. Take $M$ small enough that $LM^{2/7}<1/2$. There is an
explicit Borel policy
$D_{M,c}$ with exact budget $M$ such that
\begin{equation}\label{eq:oracle-uniform-trial}
 J_c(D_{M,c})\le C
 \begin{cases}
 M^{-1/5}t^{-4/5},&t\ge LM^{2/7},\\
 M^{-3/7},&|t|\le LM^{2/7},\\
 |t|^{-3/2},&t\le-LM^{2/7}.
 \end{cases}
\end{equation}
At each fixed $|c|<1$ it also satisfies the sharp limit
\begin{equation}\label{eq:oracle-pointwise-trial}
 M^{1/5}J_c(D_{M,c})\longrightarrow
                  2^{6/5}C_{\rm pl}(1-c^2)^{-4/5}.
\end{equation}

\paragraph{Separated roots.}
For $t\ge LM^{2/7}$ put $a=1-c^2$, let $r_1,r_2$ be the two roots,
and set
\begin{equation}\label{eq:oracle-sharp-policy}
 z=\frac{M}{a^{7/2}},\qquad \eta=z^{1/10},\qquad
 b=a(1-\eta),\qquad
 \delta=\frac{\eta\sqrt a}{2},\qquad
 R=\left(\frac{40M}{3b}\right)^{1/5}.
\end{equation}
Choosing $L$ large makes $\eta\le1/2$ and $R\le\delta/2$ throughout
this regime: $z\le L^{-7/2}$,
$R/\delta=2(40/3)^{1/5}z^{1/10}(1-\eta)^{-1/5}$.
The $\delta$-neighborhoods of the roots are disjoint. Taylor's formula
with $|\cos''|\le1$ gives
\[
 |\cos(r_j+x)-\cos r_j|\ge\sqrt a|x|-x^2/2,
 \qquad k_c(r_j+x)\ge b x^2\quad(|x|\le\delta).
\]
Place \eqref{eq:placement-profile} with curvature $b$ and mass $M/2$
on each root, and set mobility to zero elsewhere. The two masses add
exactly to $M$.

For completeness, the reciprocal-potential estimate used outside
these supports is
\begin{equation}\label{eq:oracle-reciprocal-tail}
 \int_{\Gamma\setminus\cup_j\{|s-r_j|<R\}}\frac{\dd s}{k_c(s)}
 \le \frac{C}{aR}\qquad(0<R\le c_1\sqrt a),
\end{equation}
with a fixed sufficiently small $c_1$.
Near each root, within a fixed small multiple of $\sqrt a$, Taylor's
bound gives $k_c\ge c a|s-r_j|^2$, whose truncated reciprocal integral
is $O((aR)^{-1})$. When $c$ is near a fold, write $x$ for arclength
from that fold and $z_0=2\sin(x/2)$. Then
$k_c=(z_0^2/2-t)^2$, with bounded metric factors on a fixed fold
neighborhood. Outside the two root neighborhoods, rescaling
$z_0=\sqrt{2t}\,y$ bounds the remaining reciprocal integral by
$Ct^{-3/2}\int_{\operatorname{dist}(y,\{-1,1\})>c_2}
 (y^2-1)^{-2}\dd y=O(a^{-3/2})$.
The part away from this fold is bounded. For offsets in a compact
subset of $(-1,1)$ the same estimate follows from a fixed positive
potential away from the root neighborhoods. Finally
$a^{-3/2}\le C/(aR)$, proving \eqref{eq:oracle-reciprocal-tail}.
The support bound \eqref{eq:placement-support-upper} and this estimate
prove the first line of \eqref{eq:oracle-uniform-trial}.

At fixed $|c|<1$, $\eta\to0$, $R/\delta\to0$, and $b/a\to1$.
Within each $\delta$-neighborhood the sharper exterior estimate is
$2/(bR)$ per root, while the rest is $O_c(\delta^{-1})=o(M^{-1/5})$, since
$\delta\asymp_c M^{1/10}$.
Together with the interior value $10/(bR)$ this gives the sharp
upper bound in \eqref{eq:oracle-pointwise-trial}; the compact-wall
lower theorem gives its matching lower bound.

\paragraph{Fold layer.}
For $|t|\le LM^{2/7}$, use the fold at $s_0=0$ when $c<0$, and at
$s_0=\pi$ when $c>0$. Put $r=M^{1/7}$ and choose a fixed $B$ with
$B^2/2>L+2$. The policy is
\begin{equation}\label{eq:oracle-fold-policy}
 D_{M,c}(s)=\frac{M}{2Br}\ind_{\{|s-s_0|<Br\}},
\end{equation}
where the distance is the local periodic arclength. This has mass $M$.
With $x=(s-s_0)/r$ and $\mu=t/r^2$, the potential on this fixed
interval is $r^4V_{r,\mu}(x)$, where
\[
 V_{r,\mu}(x)=\left(\frac{1-\cos(rx)}{r^2}-\mu\right)^2
 \longrightarrow (x^2/2-\mu)^2
\]
uniformly on $[-B,B]\times[-L,L]$. Since $D=r^6/(2B)$ on the patch,
the diffusion operator has the same factor $r^4$. For an unamplified
rescaled test $v(s)=\phi(x)$, the integrated form and source are
\[
 r^5\int_{-B}^B\left(\frac{|\phi'|^2}{2B}
                         +V_{r,\mu}\phi^2\right)\dd x,
 \qquad r\int_{-B}^B\phi\dd x,
\]
respectively.
The free-endpoint form on $[-B,B]$ is uniformly coercive: Poincar\'e
controls the nonconstant part, and the integral of its potential is
uniformly positive because no limiting polynomial
$(x^2/2-\mu)^2$ vanishes identically. A contradiction sequence with
unit $L^2$ norm and energy tending to zero would converge to a
nonzero constant with zero potential integral. Its source response
is therefore bounded uniformly. Since the response scale is
$r^2/r^5=r^{-3}$, restricting any smooth wall test to the patch bounds
its contribution by $Cr^{-3}$.

Outside the patch $D=0$. Its boundary has a fixed scaled margin
beyond all roots. More explicitly, for $Br\le|s-s_0|\le x_0$ and
small fixed $x_0$, $1-\cos(s-s_0)-t\ge c|s-s_0|^2$ after enlarging
$B$ if necessary. Thus the reciprocal integral outside the patch is
at most $C\int_{Br}^{x_0}x^{-4}\dd x+C\le Cr^{-3}$.
This proves the middle line of \eqref{eq:oracle-uniform-trial}
for every original smooth wall test.

\paragraph{Rootless offsets and measurability.}
For $t\le-LM^{2/7}$ use the uniform field $M/(2\pi)$.
Dropping the nonnegative derivative energy gives
$J_c(D)\le\int_\Gamma k_c^{-1}$.
Near the closest fold, $k_c(s)\ge c(|t|+|s-s_0|^2)^2$;
rescaling by $|t|^{1/2}$ gives the last bound in
\eqref{eq:oracle-uniform-trial}. Offsets away from both folds have a
bounded reciprocal integral, as required by that same bound for
$-1\le t<0$.
The policy just constructed is Borel into $L^1(\Gamma)$:
$r_1(c)=\arccos(-c)$ and $r_2(c)=2\pi-r_1(c)$ vary continuously in
the separated regime, as do its positive profile widths and
amplitudes; translations and dilations of these compact continuous
profiles are continuous in $L^1$. The two fold patches and uniform
branch are Borel, and the regime boundaries are Borel sets. Assign
either adjoining prescription at a threshold, and the fold patch at
$c=\pm1$. Thus the policy is defined with exactly the required budget
at every offset, including the null exceptional points.

\begin{proof}[Proof of \cref{thm:oracle}]
The estimates just proved give the integrable, $M$-independent bound
\begin{equation}\label{eq:oracle-domination}
 M^{1/5}J_c(D_{M,c})\le C|1-|c||^{-4/5},\qquad c\ne\pm1.
\end{equation}
For the middle regime this follows from
$M^{1/5-3/7}=M^{-8/35}$ and $|t|\le LM^{2/7}$; for the rootless
regime it follows from $M^{1/5}|t|^{-7/10}\le C$ when
$|t|\ge LM^{2/7}$.
At fixed $|c|>1$, the scaled response tends to zero. Dominated
convergence, \eqref{eq:oracle-pointwise-trial}, and the density $1/4$
therefore give
\[
 \limsup_{M\to0}M^{1/5}\mathcal O(M)
 \le\frac{2^{6/5}C_{\rm pl}}4
       \int_{-1}^1(1-c^2)^{-4/5}\dd c=K_{\rm obs}.
\]

For the converse take any sequence $M_l\downarrow0$ and Borel
policies whose expected values are at most $\mathcal O(M_l)+1$.
Such policies exist by the definition of the finite infimum; no
pointwise minimizing selection is used.
At each fixed $|c|<1$, \cref{thm:compact-placement}, applied to its
two separated zeros of curvature $a=1-c^2$, bounds every member of
this sequence and implies
\[
 \liminf_{l\to\infty}M_l^{1/5}J_c(D_{M_l,c})
 \ge 2^{6/5}C_{\rm pl}(1-c^2)^{-4/5}.
\]
Fatou's lemma applied to these nonnegative, measurable responses
proves the matching lower bound after integration over $(-1,1)$.
This argument avoids any measurability assertion for a pointwise
infimum over the uncountable coefficient class.
The full-bulk statement follows from \cref{thm:optimal-transfer},
and the ratio follows from \cref{thm:blind-sharp} at $q=1$.
\end{proof}

The fold layer has width $O(M^{2/7})$ and contributes only
$O(M^{-1/7})$ to the expected trial response. It is needed for a
valid global policy, but is smaller than the mean's leading
$M^{-1/5}$ term. Exact observation changes the exponent because each
realization can receive mobility at its actual defects; the total
budget is not redistributed among realizations. The ratio
\eqref{eq:oracle-blind-ratio} tends to zero slowly and does not imply
a specified improvement at any particular finite experimental budget.
All coefficients in this section use the dimensionless normalization
\eqref{eq:units}; the dimensional response factor must be restored
before multiplication by the physical $\chi$.
